% This is the template for submission of abstracts to NetSci 2025 in Maastricht, the Netherlands.
% It is modified from NetSci2025, which was inspired by that of NetSci2017, NetSci2016 and that of STATPHSY25.
% The editor of the booklet reserves the right to modify your submission.

% To process this file run LaTeX2e

% ******** DO NOT EDIT ****************
\documentclass[10pt]{article}
% \usepackage{mathptmx}
\usepackage[margin=20mm]{geometry}
\usepackage{graphicx}
\usepackage{lmodern}
\usepackage{hyperref}
\usepackage{xcolor}

\renewcommand{\familydefault}{\sfdefault}
\pagestyle{empty}
\setlength{\parskip}{0.25\baselineskip}
\renewcommand{\title}[1]{{\noindent\large\bfseries#1\medskip\\}}
\renewcommand{\author}[2]{{\noindent #1 \medskip\\ \noindent \small #2 \medskip\\}}
% *************************************

\begin{document}

% ********** USER DEFINED *************


% Enter title here
\title{Emergent Trail Networks}
\author{
\small
% Enter author(s) here
B.Wallace\textsuperscript{1},
A.Berdahl\textsuperscript{1},
}
{
\small
% Enter affiliation(s) here
1. University of Washington, Seattle, USA\\
}



% Enter abstract here. Please keep everything within one page.
When a person or animal walk through rough terrain, they flatten it down, making it easier for subsequent travelers to walk the same way. When multiple individuals walk between multiple destinations, this positive feedback leads to the emergence of a network of trails. Little is known of the topology, dynamics, stability and other properties of such emergent trail networks.

We describe a computational modeling approach to understanding the dynamics of these networks that builds on a classic "active-walker" model [1]. Simulated agents take the least-cost path between points of interest across a continuous 2D space; as they move they decrease the "cost" of traveling over the same patch of ground. We consider the dual optimization problem of simultaneously minimizing travel costs between points of interest in the network and total amount of path improvement, as in [2]; within this framework we try to understand the efficiencies of the emergent networks, and how efficiencies vary with environmental parameters, behavioral parameters, and other network measures. We discuss implications to urban planning, transportation system design, and animal behavior. Finally, we show that when points of interest are randomly and dynamically placed, the emergent transportation network may take on novel fractal shapes.




\bigskip
{\small
\noindent[1] Helbing, Dirk, Joachim Keltsch, and Péter Molnár. “Modelling the Evolution of Human Trail Systems.” Nature 1997

\noindent[2] Perna, Andrea, and Tanya Latty. “Animal Transportation Networks.” Journal of the Royal Society 2014\\

}



\begin{figure}[!h]
  \centering
  \includegraphics[width=.95\textwidth]{netsci2025.pdf}
  \caption{(A) Snapshots of trail development between fixed points in low-cost (left column) and high-cost (right column) environments. (B) Lines show trail development over time plotted in travel-cost (average cost of travel between points) vs infrastructure (amount of trail used) space. Final networks (large dots) show an efficient tradeoff between infrastructure use and travel cost along a pareto front. (C) Walkers move between \textit{dynamically shifting} locations, producing a network that connects all of the 2-D space. }
\end{figure}





% *************************************
\end{document}